% ============================================================
% Question Bank: ch04-02 Simplifying Expressions by Combining Like Terms
% Topic Title: Simplifying Expressions by Combining Like Terms
% CCSS: 7.EE.A.1
% Questions: 30 (18 MC + 7 SA + 5 Graphical)
% ============================================================

% --- Multiple Choice Questions (18) ---

\begin{practiceQuestion}{4-2-q01}{mc}
\begin{questionText}
Which pair of terms are like terms?
\end{questionText}
\choiceA{$3x$ and $3y$}
\choiceB{$5m$ and $-2m$}
\choiceC{$4n$ and $4n^2$}
\choiceD{$7$ and $7x$}
\correctAnswer{B}
\explanation{Like terms have the same variable raised to the same power. $5m$ and $-2m$ both have the variable $m$ to the first power.}
\end{practiceQuestion}

\begin{practiceQuestion}{4-2-q02}{mc}
\begin{questionText}
Simplify $7x + 3x$.
\end{questionText}
\choiceA{$10x$}
\choiceB{$21x$}
\choiceC{$10x^2$}
\choiceD{$73x$}
\correctAnswer{A}
\explanation{$7x$ and $3x$ are like terms. Add the coefficients: $7 + 3 = 10$. Result: $10x$.}
\end{practiceQuestion}

\begin{practiceQuestion}{4-2-q03}{mc}
\begin{questionText}
Simplify $9a - 4a + 6$.
\end{questionText}
\choiceA{$11a$}
\choiceB{$5a + 6$}
\choiceC{$5a - 6$}
\choiceD{$13a - 4$}
\correctAnswer{B}
\explanation{Combine $9a - 4a = 5a$. The constant $6$ stays. Result: $5a + 6$.}
\end{practiceQuestion}

\begin{practiceQuestion}{4-2-q04}{mc}
\begin{questionText}
Simplify $3y + 8 - 5y - 2$.
\end{questionText}
\choiceA{$8y + 6$}
\choiceB{$-2y + 6$}
\choiceC{$2y - 6$}
\choiceD{$-2y + 10$}
\correctAnswer{B}
\explanation{Combine variable terms: $3y - 5y = -2y$. Combine constants: $8 - 2 = 6$. Result: $-2y + 6$.}
\end{practiceQuestion}

\begin{practiceQuestion}{4-2-q05}{mc}
\begin{questionText}
What is the coefficient of $n$ in the expression $12 - 7n + 3$?
\end{questionText}
\choiceA{$12$}
\choiceB{$7$}
\choiceC{$-7$}
\choiceD{$3$}
\correctAnswer{C}
\explanation{The coefficient is the number multiplied by the variable. In $-7n$, the coefficient is $-7$.}
\end{practiceQuestion}

\begin{practiceQuestion}{4-2-q06}{mc}
\begin{questionText}
Simplify $-4x + 9 + 7x - 3$.
\end{questionText}
\choiceA{$3x + 6$}
\choiceB{$-11x + 6$}
\choiceC{$3x - 6$}
\choiceD{$11x + 12$}
\correctAnswer{A}
\explanation{Combine variable terms: $-4x + 7x = 3x$. Combine constants: $9 - 3 = 6$. Result: $3x + 6$.}
\end{practiceQuestion}

\begin{practiceQuestion}{4-2-q07}{mc}
\begin{questionText}
Which expression is equivalent to $6m + 4 - 2m + 3$?
\end{questionText}
\choiceA{$4m + 7$}
\choiceB{$8m + 7$}
\choiceC{$4m + 1$}
\choiceD{$4m - 7$}
\correctAnswer{A}
\explanation{Combine $6m - 2m = 4m$. Combine $4 + 3 = 7$. Result: $4m + 7$.}
\end{practiceQuestion}

\begin{practiceQuestion}{4-2-q08}{mc}
\begin{questionText}
Simplify $\frac{1}{3}x + \frac{2}{3}x$.
\end{questionText}
\choiceA{$\frac{2}{9}x$}
\choiceB{$x$}
\choiceC{$\frac{3}{6}x$}
\choiceD{$\frac{1}{3}x^2$}
\correctAnswer{B}
\explanation{The coefficients are $\frac{1}{3}$ and $\frac{2}{3}$. Add: $\frac{1}{3} + \frac{2}{3} = \frac{3}{3} = 1$. Result: $1x = x$.}
\end{practiceQuestion}

\begin{practiceQuestion}{4-2-q09}{mc}
\begin{questionText}
A student simplifies $5x + 3 + 2x$ and gets $10x$. What mistake did the student make?
\end{questionText}
\choiceA{Added all the numbers together including the constant}
\choiceB{Multiplied instead of adding}
\choiceC{Forgot to include the variable}
\choiceD{Combined unlike terms: added $3$ to $7x$ as if $3$ were $3x$}
\correctAnswer{D}
\explanation{The student treated the constant $3$ as $3x$, getting $5x + 3x + 2x = 10x$. The constant $3$ is not a like term with the variable terms. The correct answer is $7x + 3$.}
\end{practiceQuestion}

\begin{practiceQuestion}{4-2-q10}{mc}
\begin{questionText}
Simplify $2.5a + 1.3 - 0.5a + 2.7$.
\end{questionText}
\choiceA{$2a + 4$}
\choiceB{$3a + 4$}
\choiceC{$2a - 4$}
\choiceD{$6$}
\correctAnswer{A}
\explanation{Combine variable terms: $2.5a - 0.5a = 2a$. Combine constants: $1.3 + 2.7 = 4$. Result: $2a + 4$.}
\end{practiceQuestion}

\begin{practiceQuestion}{4-2-q11}{mc}
\begin{questionText}
How many terms are in the simplified form of $8b - 3 + 2b + 5 - b$?
\end{questionText}
\choiceA{$5$}
\choiceB{$3$}
\choiceC{$2$}
\choiceD{$1$}
\correctAnswer{C}
\explanation{Combine variable terms: $8b + 2b - b = 9b$. Combine constants: $-3 + 5 = 2$. Simplified: $9b + 2$, which has $2$ terms.}
\end{practiceQuestion}

\begin{practiceQuestion}{4-2-q12}{mc}
\begin{questionText}
Simplify $-6y + 4y - y$.
\end{questionText}
\choiceA{$-3y$}
\choiceB{$-y$}
\choiceC{$3y$}
\choiceD{$-11y$}
\correctAnswer{A}
\explanation{$-6y + 4y = -2y$. Then $-2y - y = -3y$.}
\end{practiceQuestion}

\begin{practiceQuestion}{4-2-q13}{mc}
\begin{questionText}
Which expression is already simplified (cannot be combined further)?
\end{questionText}
\choiceA{$4x + 3x$}
\choiceB{$5a - 2a + 7$}
\choiceC{$3m + 8$}
\choiceD{$6 + 9 - 2$}
\correctAnswer{C}
\explanation{$3m + 8$ has no like terms to combine ($3m$ is a variable term and $8$ is a constant). The others can all be simplified further.}
\end{practiceQuestion}

\begin{practiceQuestion}{4-2-q14}{mc}
\begin{questionText}
Simplify $\frac{3}{4}n - \frac{1}{4}n + 5$.
\end{questionText}
\choiceA{$\frac{1}{2}n + 5$}
\choiceB{$n + 5$}
\choiceC{$\frac{3}{4}n + 5$}
\choiceD{$\frac{2}{4}n - 5$}
\correctAnswer{A}
\explanation{Combine: $\frac{3}{4}n - \frac{1}{4}n = \frac{2}{4}n = \frac{1}{2}n$. Add the constant: $\frac{1}{2}n + 5$.}
\end{practiceQuestion}

\begin{practiceQuestion}{4-2-q15}{mc}
\begin{questionText}
Simplify $10 - 3x + 5x - 4$.
\end{questionText}
\choiceA{$8x + 6$}
\choiceB{$-8x + 6$}
\choiceC{$2x + 6$}
\choiceD{$2x - 6$}
\correctAnswer{C}
\explanation{Combine variable terms: $-3x + 5x = 2x$. Combine constants: $10 - 4 = 6$. Result: $2x + 6$.}
\end{practiceQuestion}

\begin{practiceQuestion}{4-2-q16}{mc}
\begin{questionText}
Which expression is equivalent to $-x + 8 + 4x - 5$?
\end{questionText}
\choiceA{$3x + 3$}
\choiceB{$5x + 3$}
\choiceC{$-3x + 3$}
\choiceD{$3x + 13$}
\correctAnswer{A}
\explanation{Combine variable terms: $-x + 4x = 3x$. Combine constants: $8 - 5 = 3$. Result: $3x + 3$.}
\end{practiceQuestion}

\begin{practiceQuestion}{4-2-q17}{mc}
\begin{questionText}
Simplify $5c + 2d - 3c + 4d$.
\end{questionText}
\choiceA{$8cd$}
\choiceB{$2c + 6d$}
\choiceC{$8c + 6d$}
\choiceD{$2c - 2d$}
\correctAnswer{B}
\explanation{Combine $c$-terms: $5c - 3c = 2c$. Combine $d$-terms: $2d + 4d = 6d$. Result: $2c + 6d$.}
\end{practiceQuestion}

\begin{practiceQuestion}{4-2-q18}{mc}
\begin{questionText}
A student says that $4x + 4y = 8xy$. Is the student correct?
\end{questionText}
\choiceA{Yes, because $4 + 4 = 8$ and the variables combine as $xy$.}
\choiceB{No, because $4x$ and $4y$ are not like terms and cannot be combined.}
\choiceC{Yes, because the coefficients are the same.}
\choiceD{No, because the answer should be $(4x)(4y) = 16xy$.}
\correctAnswer{B}
\explanation{$4x$ and $4y$ have different variables, so they are not like terms and cannot be combined. The expression $4x + 4y$ is already simplified.}
\end{practiceQuestion}

% --- Short Answer Questions (7) ---

\begin{practiceQuestion}{4-2-q19}{sa}
\begin{questionText}
Simplify $8p - 3 + 2p - 7$.
\end{questionText}
\correctAnswer{$10p - 10$}
\explanation{Combine variable terms: $8p + 2p = 10p$. Combine constants: $-3 - 7 = -10$. Result: $10p - 10$.}
\end{practiceQuestion}

\begin{practiceQuestion}{4-2-q20}{sa}
\begin{questionText}
Simplify $\frac{1}{2}x + \frac{3}{4}x - 2$.
\end{questionText}
\correctAnswer{$\frac{5}{4}x - 2$ or $1\frac{1}{4}x - 2$}
\explanation{Find a common denominator: $\frac{2}{4}x + \frac{3}{4}x = \frac{5}{4}x$. The constant stays: $\frac{5}{4}x - 2$.}
\end{practiceQuestion}

\begin{practiceQuestion}{4-2-q21}{sa}
\begin{questionText}
Simplify $-2a + 5b + 7a - 3b$.
\end{questionText}
\correctAnswer{$5a + 2b$}
\explanation{Combine $a$-terms: $-2a + 7a = 5a$. Combine $b$-terms: $5b - 3b = 2b$. Result: $5a + 2b$.}
\end{practiceQuestion}

\begin{practiceQuestion}{4-2-q22}{sa}
\begin{questionText}
Simplify $12 - 5w + 3w - 8 + w$.
\end{questionText}
\correctAnswer{$-w + 4$ or $4 - w$}
\explanation{Combine variable terms: $-5w + 3w + w = -w$. Combine constants: $12 - 8 = 4$. Result: $-w + 4$.}
\end{practiceQuestion}

\begin{practiceQuestion}{4-2-q23}{sa}
\begin{questionText}
The perimeter of a triangle is $3x + 2 + 5x - 1 + 2x + 4$. Write the simplified expression for the perimeter.
\end{questionText}
\correctAnswer{$10x + 5$}
\explanation{Combine variable terms: $3x + 5x + 2x = 10x$. Combine constants: $2 - 1 + 4 = 5$. Perimeter: $10x + 5$.}
\end{practiceQuestion}

\begin{practiceQuestion}{4-2-q24}{sa}
\begin{questionText}
Simplify $-3.5y + 2.1 + 1.5y - 0.6$.
\end{questionText}
\correctAnswer{$-2y + 1.5$}
\explanation{Combine variable terms: $-3.5y + 1.5y = -2y$. Combine constants: $2.1 - 0.6 = 1.5$. Result: $-2y + 1.5$.}
\end{practiceQuestion}

\begin{practiceQuestion}{4-2-q25}{sa}
\begin{questionText}
Write an expression with three terms that simplifies to $4n + 7$.
\end{questionText}
\correctAnswer{Answers vary, e.g., $n + 3n + 7$ or $5n - n + 7$}
\explanation{Any expression whose like terms combine to give $4n + 7$ is correct. For example: $n + 3n + 7 = 4n + 7$.}
\end{practiceQuestion}

% --- Graphical Questions (3 GMC + 2 GSA) ---

\begin{practiceQuestion}{4-2-q26}{gmc}
\begin{questionText}
The algebra tiles below represent an expression. Which is the simplified form?

\begin{center}
\begin{tikzpicture}[scale=0.55]
  % Positive x-tiles
  \fill[funBlue!40] (0,0) rectangle (2,1);
  \draw[thick] (0,0) rectangle (2,1);
  \node at (1,0.5) {$x$};
  \fill[funBlue!40] (2.3,0) rectangle (4.3,1);
  \draw[thick] (2.3,0) rectangle (4.3,1);
  \node at (3.3,0.5) {$x$};
  \fill[funBlue!40] (4.6,0) rectangle (6.6,1);
  \draw[thick] (4.6,0) rectangle (6.6,1);
  \node at (5.6,0.5) {$x$};
  \fill[funBlue!40] (6.9,0) rectangle (8.9,1);
  \draw[thick] (6.9,0) rectangle (8.9,1);
  \node at (7.9,0.5) {$x$};
  \fill[funBlue!40] (9.2,0) rectangle (11.2,1);
  \draw[thick] (9.2,0) rectangle (11.2,1);
  \node at (10.2,0.5) {$x$};
  % Negative x-tiles
  \fill[funRed!40] (0,-1.5) rectangle (2,-0.5);
  \draw[thick] (0,-1.5) rectangle (2,-0.5);
  \node at (1,-1) {$-x$};
  \fill[funRed!40] (2.3,-1.5) rectangle (4.3,-0.5);
  \draw[thick] (2.3,-1.5) rectangle (4.3,-0.5);
  \node at (3.3,-1) {$-x$};
  % Positive unit tiles
  \fill[funGreen!40] (6,-1.5) rectangle (7,-0.5);
  \draw[thick] (6,-1.5) rectangle (7,-0.5);
  \node at (6.5,-1) {$1$};
  \fill[funGreen!40] (7.3,-1.5) rectangle (8.3,-0.5);
  \draw[thick] (7.3,-1.5) rectangle (8.3,-0.5);
  \node at (7.8,-1) {$1$};
  \fill[funGreen!40] (8.6,-1.5) rectangle (9.6,-0.5);
  \draw[thick] (8.6,-1.5) rectangle (9.6,-0.5);
  \node at (9.1,-1) {$1$};
  \fill[funGreen!40] (9.9,-1.5) rectangle (10.9,-0.5);
  \draw[thick] (9.9,-1.5) rectangle (10.9,-0.5);
  \node at (10.4,-1) {$1$};
\end{tikzpicture}
\end{center}
\end{questionText}
\choiceA{$5x + 4$}
\choiceB{$3x + 4$}
\choiceC{$7x + 4$}
\choiceD{$3x - 4$}
\correctAnswer{B}
\explanation{There are $5$ positive $x$-tiles and $2$ negative $x$-tiles: $5x + (-2x) = 3x$. There are $4$ positive unit tiles: $+4$. Simplified: $3x + 4$.}
\end{practiceQuestion}

\begin{practiceQuestion}{4-2-q27}{gmc}
\begin{questionText}
The lengths of the four sides of a quadrilateral are shown below. What is the simplified perimeter?

\begin{center}
\begin{tikzpicture}[scale=0.9]
  \draw[thick, fill=funBlue!10] (0,0) -- (4,0) -- (5,3) -- (1,3) -- cycle;
  \node[below] at (2,0) {$3x + 2$};
  \node[right] at (4.7,1.5) {$x + 1$};
  \node[above] at (3,3) {$2x - 3$};
  \node[left] at (0.3,1.5) {$x + 4$};
\end{tikzpicture}
\end{center}
\end{questionText}
\choiceA{$7x + 4$}
\choiceB{$6x + 4$}
\choiceC{$7x - 4$}
\choiceD{$6x + 10$}
\correctAnswer{A}
\explanation{Add all sides: $(3x + 2) + (x + 1) + (2x - 3) + (x + 4)$. Combine variable terms: $3x + x + 2x + x = 7x$. Combine constants: $2 + 1 - 3 + 4 = 4$. Perimeter: $7x + 4$.}
\end{practiceQuestion}

\begin{practiceQuestion}{4-2-q28}{gmc}
\begin{questionText}
Look at the two groups of algebra tiles below. Which expression represents their sum?

\begin{center}
\begin{tikzpicture}[scale=0.55]
  % Group 1 label
  \node[font=\bfseries] at (3, 1.5) {Group 1};
  \fill[funBlue!40] (0,0) rectangle (2,1);
  \draw[thick] (0,0) rectangle (2,1);
  \node at (1,0.5) {$x$};
  \fill[funBlue!40] (2.3,0) rectangle (4.3,1);
  \draw[thick] (2.3,0) rectangle (4.3,1);
  \node at (3.3,0.5) {$x$};
  \fill[funGreen!40] (4.8,0) rectangle (5.8,1);
  \draw[thick] (4.8,0) rectangle (5.8,1);
  \node at (5.3,0.5) {$1$};

  % Group 2 label
  \node[font=\bfseries] at (3, -1) {Group 2};
  \fill[funBlue!40] (0,-2.5) rectangle (2,-1.5);
  \draw[thick] (0,-2.5) rectangle (2,-1.5);
  \node at (1,-2) {$x$};
  \fill[funGreen!40] (2.5,-2.5) rectangle (3.5,-1.5);
  \draw[thick] (2.5,-2.5) rectangle (3.5,-1.5);
  \node at (3,-2) {$1$};
  \fill[funGreen!40] (3.8,-2.5) rectangle (4.8,-1.5);
  \draw[thick] (3.8,-2.5) rectangle (4.8,-1.5);
  \node at (4.3,-2) {$1$};
  \fill[funGreen!40] (5.1,-2.5) rectangle (6.1,-1.5);
  \draw[thick] (5.1,-2.5) rectangle (6.1,-1.5);
  \node at (5.6,-2) {$1$};
\end{tikzpicture}
\end{center}
\end{questionText}
\choiceA{$3x + 3$}
\choiceB{$2x + 4$}
\choiceC{$3x + 4$}
\choiceD{$2x + 3$}
\correctAnswer{C}
\explanation{Group 1: $2x + 1$. Group 2: $x + 3$. Sum: $(2x + 1) + (x + 3) = 3x + 4$.}
\end{practiceQuestion}

\begin{practiceQuestion}{4-2-q29}{gsa}
\begin{questionText}
A rectangular garden has the dimensions shown. Write a simplified expression for the total distance around the garden (its perimeter).

\begin{center}
\begin{tikzpicture}[scale=0.7]
  \draw[thick, fill=funGreen!10] (0,0) rectangle (6,3);
  \node[below] at (3,0) {$4x - 1$};
  \node[right] at (6,1.5) {$2x + 5$};
  \draw[thick, <->] (0,-0.7) -- (6,-0.7);
  \draw[thick, <->] (6.7,0) -- (6.7,3);
\end{tikzpicture}
\end{center}
\end{questionText}
\correctAnswer{$12x + 8$}
\explanation{Perimeter $= 2(4x - 1) + 2(2x + 5) = 8x - 2 + 4x + 10 = 12x + 8$.}
\end{practiceQuestion}

\begin{practiceQuestion}{4-2-q30}{gsa}
\begin{questionText}
The table shows the number of each type of tile in a bag. Write a simplified expression for the total number of tiles.

\begin{center}
\begin{tabular}{|l|c|}
\hline
\textbf{Tile Color} & \textbf{Number of Tiles} \\
\hline
Red & $3n + 4$ \\
\hline
Blue & $2n - 1$ \\
\hline
Green & $n + 5$ \\
\hline
Yellow & $4n - 2$ \\
\hline
\end{tabular}
\end{center}
\end{questionText}
\correctAnswer{$10n + 6$}
\explanation{Add all expressions: $(3n + 4) + (2n - 1) + (n + 5) + (4n - 2)$. Combine variable terms: $3n + 2n + n + 4n = 10n$. Combine constants: $4 - 1 + 5 - 2 = 6$. Total: $10n + 6$.}
\end{practiceQuestion}
